\subsection{The ordinary projector corner}
For a rank-$j$ projector in dimension $h$ with $j>h/2$, put $r=h-j$.
In a generic rational basis its first-$r$/last-$r$ corner is invertible:
one witness pairs these coordinates in rank-one blocks
$\frac12\left(\begin{smallmatrix}1&1\\1&1\end{smallmatrix}\right)$
and is the identity on the middle $h-2r$ coordinates. Similarity of
rational projectors with the same rank proves that the relevant minor
is a nonzero rational function of the basis.

Write $I-P=XY$, $YX=I_r$. The corner is $-X_LY_R$, so $X_L,Y_R$ are
invertible. Eliminating it leaves on the middle coordinates
\[
 I-X_MY_M-(-X_MY_R)(-Y_R^{-1}X_L^{-1})(-X_LY_M)=I_{h-2r}.
\]
The resulting lower-lower profile is $r$ singleton pivots and one
contiguous block of width $2j-h$. For $j\le h/2$ retain $j$ singleton
pivots; for $j=h$ use one block of width $h$. The last is still smaller
than the full tensor dimension $m$. Nonzero scalings are absorbed in
the triangular factors. This argument is also the large-projector
identity used in the controlled corners below.

\subsection{Inner boundary prescriptions}
Use $F_1=\mathbb Q^{25}$, $F_2=\mathbb Q^{23}$, $F_3=\mathbb Q^{57}$,
$H=575$ and $m=32775$. For physical inner index $i=23\alpha+\beta$ set
\[
 K=T_2(I_{25}\otimes L_2),\qquad
 T_2(u\otimes e_\beta)=J_\beta u\otimes e_\beta,
 \qquad J_\beta=M_\beta L_1,
\]
where $L_1,L_2$ are free invertible rational bases. Put $n=82$,
$I=(0,\ldots,24)$, $\Sigma=(0,\ldots,6)$ and
\[
 R=[I,\Sigma,I,I],\qquad C=[I,I,\Sigma,I].
\]
For $i<n$ prescribe $e_\alpha^tM_\beta=e_{R_i}^t$; for $i\ge H-n$
prescribe $M_\beta^{-1}e_\alpha=e_{C_{i-(H-n)}}$.
The latter is equivalently a fixed column of $M_\beta$.
The effective starting residues modulo 23 are respectively
$\{0,2,9,11\}$ and $\{10,12,14,21\}$.
They are pairwise distinct and disjoint. No coefficient coordinate is
prescribed twice on one side, nor shared across the two sides. The
prescribed row positions are near the start and the prescribed column
positions near the end. They are disjoint, so every $M_\beta$ completes
to a permutation matrix. Their unspecified entries form an affine
space; its invertible locus is nonempty and irreducible.

The first 25 rows and last 25 inverse columns give matching $L_1$ and
$L_1^{-1}$ bases, so every stage-one single-active-factor projector has
an ordinary 25-corner. For stage two, the first factor is rank one;
its row/column weights separate, leaving a diagonal rescaling of
$L_2BL_2^{-1}$ in the first-23/last-23 corner. Thus all ordinary
23-dimensional profiles are available in the same family.

\subsection{The A5 tree certificate}
The inner projector $(I-P_1)\otimes(I-P_2)$ has rank $528$ and nullity
$47$. Eliminate the first 23 rows of its nullspace restriction. The
remaining 24 weighted incidence vectors on $\{0,\ldots,24\}$ have
forward edges
\[
 (0,23),(1,24);\quad ((j+2)\bmod23,j)\ (0\le j<7);\quad
 ((j+9)\bmod23,j)\ (0\le j\le14)
\]
and dual edges
\[
 (23,0),(24,1);\quad (2+((-9+j)\bmod23),j)\ (0\le j<7);\quad
 (2+((-11+j)\bmod23),j)\ (10\le j\le24).
\]
Both are trees. Here are leaf-removal orders, each leaving vertex 24:
\begin{align*}
 \text{forward: }&15,16,7,17,8,6,18,9,19,10,20,11,21,12,22,13,\\
 &4,2,23,0,14,5,3,1;\\
 \text{dual: }&7,8,9,16,0,17,1,18,2,11,20,4,13,22,6,15,\\
 &23,14,5,21,12,3,19,10.
\end{align*}
For nonzero line coordinates, diagonal rescaling converts these into
ordinary tree-incidence vectors. Leaf elimination proves rank 24.
Both 47-dimensional nullspace restrictions are therefore nonsingular;
the large-projector identity gives 47 singleton pivots and one
contiguous block $575-94=481$.

\subsection{Matched outer basis and the paired A1 corner}
Choose
\[
 S=T(I_{57}\otimes K),\qquad T(u\otimes e_i)=G_i u\otimes e_i,
\]
with a free $L_0\in\operatorname{GL}_{57}(\mathbb Q)$.
Prescribe the first 57 outer row functionals in the $L_0$ coordinate
order and the next 25 as repetitions of its first 25 coordinates.
Prescribe the last 57 inverse columns in matching $L_0^{-1}$ order
and the preceding 25 as repetitions of its first 25 coordinates.
The first and last 82 multiplicity intervals are disjoint and act on
disjoint $G_i$, giving nonempty rational completion families. Other
entries remain free. In particular the ordinary 57-dimensional
stage-three corners and the full 57-block for an unmatched entrance
are retained.

For a paired role the null projector is
\[
 Q=I_{57}\otimes P_q+P_t\otimes P_U,\qquad
 \operatorname{rank}P_q=1,\quad \operatorname{rank}P_U=25,\quad
 P_qP_U=P_UP_q=0.
\]
It has rank 82 and the actual A1 projector is $I-Q$.
Factor $Q=XY$, $YX=I_{82}$, and consider the first-82/last-82 corner.
In the first 25 rows against the rightmost 25 columns, the outer
coordinate groups are disjoint, so $I_{57}\otimes P_q$ contributes
zero. The $P_t\otimes P_U$ contribution is a generically nonzero
diagonal. Normalize and eliminate this contiguous 25-block.
These pivots exhaust its 25-dimensional coefficient contribution.
The next seven rows have a diagonal seven-block and entries only to
its left, supplying the next contiguous block. The residual 50-square
corner is a two-by-two matrix of diagonal 25-blocks: its top-right
diagonal is generically nonzero, and its determinant is nonzero, so it
supplies two more contiguous 25-blocks.

Here is the nonvanishing argument in the prescribed parameter family.
The duplicated outer groups in $X$ compare primal coefficient ratios
at multiplicity shift $57\bmod23=11$; those in $Y$ compare dual ratios
at shift $25\bmod23=2$. Distinct coefficient positions are free and
the ratios are generically distinct. Orthogonality of the two motif
lines imposes their bilinear pairing, not equality of these distinct
coordinate ratios. Thus both restrictions of $X,Y$ are invertible.
After the first block elimination, the second diagonal pivot has a
reciprocal-first-pivot term with nonzero coefficient in the free outer
rank-one coordinates, so it is not identically zero. The last pivot
then follows from full nonsingularity. The large-projector identity
adds the untouched middle block of width $m-164=32611$. The complete
ordered A1 profile is consequently
\[
 25,7,25,25,32611,
 \qquad 25+7+25+25+32611=m-82.
\]

\subsection{Remaining macro corners and simultaneous choice}
A2 is $(I-P_t)\otimes I_H$: its first-$H$/last-$H$ corner is diagonal,
giving $575,31625$. A4 has an identity second inner factor, so its
first-23/last-23 corner is diagonal with generically nonzero entries,
giving $23,529$.

For A3 the top-right 57-square corner within the $H$-coordinate block
is a diagonal rescaling of an ordinary rank-56 projector. Its first
and last 57 multiplicity intervals are disjoint; the diagonal tensor
contribution therefore vanishes there. It supplies a singleton and
a 55-block. After these exhaust the rank-56 perturbation, the untouched
matching interval supplies $H-114+2=463$ diagonal pivots as one block.
The remaining finite nullity corner supplies 112 singleton pivots;
the large-projector middle block has width $m-2H-114+2=31513$.
This gives 113 singleton pivots and blocks $55,463,31513$.

Every additional required minor is a nonzero polynomial after clearing
determinants from denominators: ordinary corners have the explicit
projector witnesses above; diagonal corners use independent nonzero
line coordinates; A5 uses the two trees; A1 uses the ratio and pivot
argument. There are only finitely many producer edges. Their product
with all invertibility determinants is nonzero on the same irreducible
rational parameter family. A rational point avoiding it exists and can
be found by enumerating rational parameters and testing these minors.
This is fixed-machine preparation independent of recursive widths.

\subsection{Complete child list and exact bit moment}
With $B_i,E,N$ as above, the macro classes are
\[
\begin{array}{c|r|l}
 \text{class}&\text{multiplicity}&\text{child widths}\\\hline
 A1&B_1&25,7,25,25,32611\\
 \text{unpaired}&E&57,32661\\
 A2&B_2&575,31625\\
 A4&B_2&23,529\\
 A3&2N&113\text{ singletons};55,463,31513\\
 A5&2N&47\text{ singletons};481
\end{array}
\]
Physical data growth adds $2N$ copies of each pair $(1,23)$, $(1,21)$,
$(1,55)$. For each $h\in\{25,23,57\}$ and each histogram rank $j$,
replicate its ordinary profile $NH_{h,j}/v_h$ times. If $j>h/2$ it is
$h-j$ singletons and one width-$(2j-h)$ block; otherwise it is $j$
singletons. Width-one blocks are folded into
$c_1=179603173231020$. These prescriptions define every $c_t$ exactly,
without a separate choice of producer histogram. Their rank sum is
\[
 \sum_t c_tt=s=Wm-2N+2L.
\]
All widths are strictly less than $m$. Define
\[
 \Psi(x)=\frac{\sum_tc_tt^x}{Wm^x},\qquad
 a_b=\frac{188459}{50000000000},\quad \tau=1-a_b.
\]
For completeness the exact rational enclosure used for the certificate
is as follows. For $1\le y\le2$, put $z=(y-1)/(y+1)$ and
\[
 L(y)=2\sum_{j=0}^{23}\frac{z^{2j+1}}{2j+1}
       +\frac{2z^{49}}{49(1-z^2)}\ \ge\log y.
\]
Extract powers of two in $m/t$ and round the resulting upper logarithm
up to a multiple of $10^{-12}$, calling it $\ell_t$.
Since $e^u\le1+u+u^2/(2(1-u/3))$ for $0\le u<3$,
\[
 \Psi(\tau)\le\sum_t\frac{c_tt}{Wm}
 \left(1+a_b\ell_t+\frac{(a_b\ell_t)^2}{2(1-a_b\ell_t/3)}\right)
 <1-6.1778\times10^{-14}.
\]
The first inequality follows termwise from
$(t/m)^{1-a_b}=(t/m)\exp(a_b\log(m/t))$; the final comparison is an
exact rational certificate. Together with the integer-width and row
construction below, it supplies the retained interchange interface with
exponent $\tau$.
